\section{Research Questions}

This thesis investigates how heterogeneous subsystem models from different simulation environments can be integrated into a coherent system-level simulation workflow.
The central research question is formulated as follows:

\textit{How can a framework support heterogeneous hybrid co-simulation of subsystem models from different tools while enabling consistent orchestration, event handling, and runtime switching between alternative models of the same subsystem?}

To answer this main question, the following sub-questions are addressed:

\begin{enumerate}
    \item[\textbf{RQ1}] How can heterogeneous simulation models from equation-based, musculoskeletal, and finite element tools be represented through a unified component interface for co-simulation?

    \item[\textbf{RQ2}] How can dependencies between heterogeneous components be analyzed in order to derive a consistent execution order and to identify direct feedthrough and algebraic loops?

    \item[\textbf{RQ3}] How can hybrid phenomena, such as contact, switching, and discrete state changes, be detected, localized, and handled within the same co-simulation workflow?

    \item[\textbf{RQ4}] How can runtime switching between alternative subsystem models, including low-dimensional and more detailed representations, be achieved with sufficiently consistent state transfer?

    \item[\textbf{RQ5}] To what extent can the resulting framework reproduce relevant system behavior and support representative benchmark scenarios such as the controlled pendulum?
\end{enumerate}


